\documentclass[12pt]{article}
\usepackage[utf8]{inputenc}
\usepackage[T1]{fontenc}
\usepackage[margin=1in]{geometry}
\usepackage{setspace}
\usepackage{xcolor}
\usepackage[colorlinks=true, urlcolor=blue!60!black, linkcolor=blue!60!black]{hyperref}
\setlength{\parindent}{2em}
\setlength{\parskip}{1.25em}
\renewcommand{\baselinestretch}{1.0}
\title{\textit{Moore v. United States}, 602 U.S. 572 (2024)}
\author{Case Brief}
\date{}
\begin{document}
\maketitle

\section{Case Metadata}

\begin{itemize}
  \item \textbf{Case Name:} \textit{Moore v. United States} (styled \textit{Charles G. Moore, et ux. v. United States})
  \item \textbf{Citation:} 602 U.S. 572 (2024); Docket No.\ 22-800
  \item \textbf{Briefing Book Placement:} Chapter 3: Federal Legislative Power (Taxing Power; Direct Tax Clauses and Sixteenth Amendment)
  \item \textbf{Dates:} Certiorari Granted June 26, 2023; Argued December 5, 2023; Decided June 20, 2024
  \item \textbf{Vote:} 7--2 (affirming the U.S. Court of Appeals for the Ninth Circuit, 36 F.4th 930 (2022))
  \item \textbf{Opinion of the Court:} Justice Kavanaugh, joined by Chief Justice Roberts and Justices Sotomayor, Kagan, and Jackson
  \item \textbf{Concurring Opinion:} Justice Jackson (joining the Court's opinion in full)
  \item \textbf{Concurrence in the Judgment:} Justice Barrett, joined by Justice Alito
  \item \textbf{Dissenting Opinion:} Justice Thomas, joined by Justice Gorsuch
  \item \textbf{Primary Sources:} \href{https://www.supremecourt.gov/opinions/23pdf/602us1r41_h3dj.pdf}{SupremeCourt.gov Official Page Proof}; \href{https://supreme.justia.com/cases/federal/us/602/22-800/}{Justia U.S. Supreme Court Center}; \href{https://www.law.cornell.edu/supremecourt/text/22-800_2024-06-20}{Cornell LII}
\end{itemize}

\section{Facts and Procedural History}

In 2006, Charles and Kathleen Moore invested \$40,000 in KisanKraft Machine Tools Private Ltd.\ (KisanKraft), an American-controlled foreign corporation (CFC) founded by a friend in India to supply equipment to small farmers, receiving a 13\% equity stake (\href{https://supreme.justia.com/cases/federal/us/602/22-800/}{Justia}). Between 2006 and 2017, KisanKraft generated substantial profits every year and reinvested them into the business rather than distributing dividends. Consequently, neither KisanKraft nor the Moores paid U.S.\ federal income tax on those accumulated earnings.

Under pre-2017 law, Subpart F of the Internal Revenue Code (26 U.S.C.\ \S\S\ 951--952) taxed U.S.\ shareholders of CFCs on a pass-through basis primarily on passive foreign income, allowing trillions of dollars in active foreign corporate earnings to remain untaxed offshore. In the Tax Cuts and Jobs Act of 2017 (TCJA), 131 Stat.\ 2054, Congress transitioned toward a more territorial international corporate tax regime and enacted the Mandatory Repatriation Tax (MRT), 26 U.S.C.\ \S\ 965. The MRT imposed a one-time, backward-looking pass-through tax (at rates of 8\% to 15.5\%) on U.S.\ shareholders holding at least a 10\% stake in a CFC, attributing their pro rata shares of the CFC's accumulated post-1986 undistributed earnings (\href{https://www.law.cornell.edu/supremecourt/text/22-800_2024-06-20}{Cornell LII}).

By the end of the 2017 tax year, the Moores' pro rata share of KisanKraft's accumulated earnings totaled roughly \$508,000. After an applicable deduction, they declared \$132,512 in MRT income and paid \$14,729 in tax. They then sued for a refund in the U.S. District Court for the Western District of Washington (No.\ C19-1539-JCC, 2020 WL 6799022), claiming that (1) the MRT was an unapportioned direct tax on property violating Article I, \S\ 2, cl.\ 3, Article I, \S\ 9, cl.\ 4, and the Sixteenth Amendment, and (2) its retroactive reach violated Fifth Amendment Due Process. The district court dismissed the complaint under Rule 12(b)(6). The Ninth Circuit affirmed, 36 F.4th 930 (9th Cir.\ 2022), declaring that ``realization of income is not a constitutional requirement'' and rejecting the Due Process claim. The Ninth Circuit denied rehearing en banc, 53 F.4th 507 (9th Cir.\ 2022), over a four-judge dissent by Judge Bumatay. The Moores sought certiorari solely on the Direct Tax Clause and Sixteenth Amendment issue; the Supreme Court granted review (599 U.S.\ \_\_\_ (2023)) after Justice Alito issued a September 8, 2023 statement declining calls to recuse over \textit{Wall Street Journal} interviews co-authored by petitioners' counsel (\href{https://supreme.justia.com/cases/federal/us/602/22-800/}{Justia}).

\section{Issue(s) Presented}

\begin{enumerate}
  \item \textbf{Question Framed by Petitioners:} Whether the Sixteenth Amendment authorizes Congress to tax unrealized sums without apportionment among the States under Article I, \S\ 2, cl.\ 3 and \S\ 9, cl.\ 4.
  \item \textbf{Narrower Question Decided by the Court:} Whether Congress may, consistent with Article I, \S\S\ 8--9 and the Sixteenth Amendment, attribute an entity's realized and undistributed income to its shareholders or partners and tax them on their pro rata portions of that income without apportionment when the entity itself is not taxed on that income.
\end{enumerate}

\section{Holding}

Yes. In a 7--2 decision authored by Justice Kavanaugh, the Court affirmed the Ninth Circuit and held that the 2017 Mandatory Repatriation Tax falls squarely within Congress's constitutional taxing power (\href{https://www.law.cornell.edu/supremecourt/text/22-800_2024-06-20}{Cornell LII}). When an entity such as a partnership, S corporation, or controlled foreign corporation realizes income that remains undistributed and untaxed at the entity level, Congress may attribute that realized income to the entity's shareholders or partners and tax them on their pro rata shares without apportionment. The Court expressly stressed that its holding is narrow: it applies only to pass-through taxation and explicitly reserves whether Congress may (1) tax both an entity and its owners on the same undistributed income, (2) enact an unapportioned tax on wealth, holdings, or net worth, or (3) tax unrealized appreciation without a constitutional realization event.

\section{Reasoning}

\subsection{Constitutional Architecture of the Taxing Power}
Justice Kavanaugh began with Article I, \S\ 8, cl.\ 1, which grants Congress an ``exhaustive'' power to lay and collect two classes of taxes: \textit{direct taxes} (generally on persons or property, requiring politically unworkable population-based apportionment under Article I, \S\ 2, cl.\ 3 and \S\ 9, cl.\ 4) and \textit{indirect taxes} (on activities or transactions, requiring only geographical uniformity). Before 1895, federal income taxes were understood as indirect taxes (\textit{Springer v. United States}, 102 U.S.\ 586 (1881)). After \textit{Pollock v. Farmers' Loan \& Trust Co.}, 158 U.S.\ 601 (1895), treated taxes on income from property as direct taxes on the underlying property, the 1913 Sixteenth Amendment repudiated \textit{Pollock}'s conflation and restored taxes on ``incomes, from whatever source derived,'' to the non-apportioned indirect class (\textit{Brushaber v. Union Pacific R. Co.}, 240 U.S.\ 1 (1916)) (\href{https://www.supremecourt.gov/opinions/23pdf/602us1r41_h3dj.pdf}{SupremeCourt.gov}).

\subsection{Entity Realization and Attribution Precedents}
Rather than resolving whether the Sixteenth Amendment mandates realization, the Court held that the MRT \textit{does} tax realized income---namely, profits realized by KisanKraft and attributed to its U.S.\ shareholders. A settled line of post-Sixteenth Amendment precedents establishes that Congress may tax either a business entity or its owners on the entity's undistributed income: \textit{Burk-Waggoner Oil Ass'n v. Hopkins}, 269 U.S.\ 110 (1925) (state-law entity labels do not constrain Congress's choice of taxpayer); \textit{Burnet v. Leininger}, 285 U.S.\ 136 (1932) (Congress may tax a partnership or its partners); \textit{Heiner v. Mellon}, 304 U.S.\ 271 (1938) (partners taxable on undistributed partnership profits even when state law barred distribution); and \textit{Helvering v. National Grocery Co.}, 304 U.S.\ 282 (1938) (extending \textit{Heiner}'s attribution principle to corporate shareholders). The Court dismissed the Moores' reliance on \textit{Eisner v. Macomber}, 252 U.S.\ 189 (1920), noting that \textit{Macomber} invalidated a tax on pro rata stock dividends where neither corporation nor shareholder realized new gain, and its remark that accumulated corporate profits are ``capital, not income'' was dicta not addressing pass-through attribution (\href{https://supreme.justia.com/cases/federal/us/602/22-800/}{Justia}).

\subsection{Longstanding Congressional Practice and Fiscal Consequences}
Reinforcing judicial precedent, Congress has taxed owners on undistributed entity profits since the 1864 Civil War income tax (upheld in \textit{Collector v. Hubbard}, 12 Wall.\ 1 (1871)), continuing through the 1913 partnership tax, 1918 personal service corporation tax, 1937 foreign personal holding company tax, 1958 S-corporation regime, and 1962 Subpart F. The Court rejected the Moores' three ``ad hoc distinctions'' attempting to preserve those concededly valid taxes while invalidating only the MRT: (1) partnerships were treated as separate taxable entities in many states and under the 1917 Revenue Act; (2) S-corporation pass-through status cannot rest on ``consent'' because 49\% minority shareholders cannot revoke it and private consent cannot waive a structural constitutional limit; and (3) ``constructive realization'' is an invented label that cannot separate Subpart F from the MRT, as both require a 10\% ownership threshold and deter offshore deferral. Taken to its logical conclusion, the Moores' theory would jeopardize vast swaths of the Internal Revenue Code (partnerships, S corporations, Subpart F, GILTI under \S\ 951A, original-issue discount under \S\ 1272, and mark-to-market futures under \S\ 1256), causing trillions in lost revenue---a ``fiscal calamity'' the Constitution does not require. Finally, the Court noted that Due Process bars arbitrary attribution (\textit{Burnet v. Wells}, 289 U.S.\ 670 (1933)), while reserving wealth, net-worth, and unrealized-appreciation taxes ``for another day'' (\href{https://www.law.cornell.edu/supremecourt/text/22-800_2024-06-20}{Cornell LII}).

\section{Concurring and Dissenting Opinions}

\begin{itemize}
  \item \textbf{Justice Jackson (Concurring):} Joining the majority in full, Justice Jackson emphasized two additional hurdles facing taxpayers who challenge uniform federal taxes. First, the Sixteenth Amendment's text contains no realization requirement; \textit{Macomber}'s Court-created limit was subsequently confined to its facts (\textit{Helvering v. Bruun}, 309 U.S.\ 461 (1940); \textit{Cottage Savings Ass'n v. Commissioner}, 499 U.S.\ 554 (1991)). Second, even outside the Sixteenth Amendment, a tax on doing business through a CFC could independently be sustained as an indirect \textit{excise} under Article I, \S\ 8 (\textit{Flint v. Stone Tracy Co.}, 220 U.S.\ 107 (1911)). Citing Justice Harlan's \textit{Pollock} dissent, she urged judicial restraint: remedies for unpopular taxes belong at the ballot box (\href{https://supreme.justia.com/cases/federal/us/602/22-800/}{Justia}).
  \item \textbf{Justice Barrett (Concurring in the Judgment, joined by Justice Alito):} Justice Barrett directly answered the question presented in the negative: Sixteenth Amendment ``income'' requires realization---converting or deriving gain from capital---and the Moores personally realized nothing from KisanKraft. She criticized the majority for misreading \textit{Burk-Waggoner}, \textit{Leininger}, \textit{Heiner}, and \textit{National Grocery} to grant Congress broad power to disregard corporate form, warning that taxing shareholders of a widely held domestic corporation would raise serious constitutional problems. She concurred in the judgment solely because the Moores conceded Subpart F's constitutionality and failed to carry their refund burden to distinguish the MRT from Subpart F (\href{https://www.law.cornell.edu/supremecourt/text/22-800_2024-06-20}{Cornell LII}).
  \item \textbf{Justice Thomas (Dissenting, joined by Justice Gorsuch):} Justice Thomas argued that the Direct Tax Clause was an essential federalism compromise protecting state tax bases, and that the Sixteenth Amendment overruled only \textit{Pollock}'s conflation of income with its source---leaving property taxes subject to apportionment. Distinguishing ``income'' from its ``source'' requires realization as defined in \textit{Macomber}. Accusing the majority of dodging the question presented by inventing an ``attribution'' doctrine from unrelated tax-evasion and statutory cases, he contended the MRT unconstitutionally taxed mere stock ownership on a fixed trigger date. He also dismissed the majority's consequentialist fear of ``fiscal calamity'' and its cautionary dicta reserving wealth taxes as ``cheap dicta'' (\href{https://supreme.justia.com/cases/federal/us/602/22-800/}{Justia}).
\end{itemize}

\section{Key Quotations}

\begin{quote}
``So the precise and narrow question that the Court addresses today is whether Congress may attribute an entity's realized and undistributed income to the entity's shareholders or partners, and then tax the shareholders or partners on their portions of that income. This Court's longstanding precedents, reflected in and reinforced by Congress's longstanding practice, establish that the answer is yes.'' --- Justice Kavanaugh, Opinion of the Court (\href{https://supreme.justia.com/cases/federal/us/602/22-800/}{Justia}).
\end{quote}

\begin{quote}
``Moreover, both before and after the Sixteenth Amendment was adopted, the term `income' was widely recognized as flexible enough to include both realized and unrealized gains.'' --- Justice Jackson, concurring (\href{https://www.law.cornell.edu/supremecourt/text/22-800_2024-06-20}{Cornell LII}).
\end{quote}

\begin{quote}
``The Sixteenth Amendment's reference to income `derived' from any source encompasses a requirement that income, to be taxed without apportionment, must be realized.'' --- Justice Barrett, concurring in the judgment (\href{https://www.law.cornell.edu/supremecourt/text/22-800_2024-06-20}{Cornell LII}).
\end{quote}

\begin{quote}
``Sixteenth Amendment `incomes' include only income realized by the taxpayer. The text and history of the Amendment make clear that it requires a distinction between `income' and the `source' from which that income is `derived.' And, the only way to draw such a distinction is with a realization requirement.'' --- Justice Thomas, dissenting (\href{https://supreme.justia.com/cases/federal/us/602/22-800/}{Justia}).
\end{quote}

\section{Doctrinal Significance and Modern Legacy (Through 2026)}

\textit{Moore} averted a seismic disruption of federal pass-through and international taxation (protecting Subpart F, GILTI, S corporations, and partnerships) while leaving the ultimate constitutional fate of federal wealth and mark-to-market taxes formally undecided (\href{https://harvardlawreview.org/print/vol-138/moore-v-united-states/}{Harvard Law Review}). Yet head-counting the opinions reveals a formidable constitutional warning shot: four Justices (Barrett, Alito, Thomas, and Gorsuch) explicitly endorsed a Sixteenth Amendment realization requirement, only Justice Jackson explicitly rejected it, and four majority Justices (Kavanaugh, Roberts, Sotomayor, and Kagan) reserved the issue while highlighting the Government's own concession that an unapportioned net-worth tax might be a direct property tax (\href{https://ir.lawnet.fordham.edu/cgi/viewcontent.cgi?article=2418&context=faculty_scholarship}{Gamage, Brooks \& McCaffery}).

Through 2026, \textit{Moore}'s doctrinal legacy has unfolded along three tracks: (1) \textbf{Federal Reform Pivots:} Because Congress's 2025 tax legislation focused on extending 2017 TCJA provisions and \textit{Moore} cast constitutional doubt over a freestanding federal wealth tax, scholars and lawmakers by 2026 pivoted toward alternative designs---such as incremental ``mix-and-match'' reforms against ``buy, borrow, die,'' proposed corporate wealth taxes framed as excises (\href{https://repository.law.umich.edu/cgi/viewcontent.cgi?article=4225&context=facarticles}{Avi-Yonah, Oct.\ 2026}), consumption-tax arguments (\href{https://washulawreview.org/consumption-taxes-and-the-constitution/}{Wash.\ U.\ L.\ Rev.}), and Senate proposals like the May 2026 Fair Trusts for Fiscal Responsibility Act (S.\ 4490), imposing an annual trust withholding tax creditable against the estate tax (\href{https://wealthcarelawyer.com/articles/briefings/who-is-writing-the-next-estate-tax-bill/}{WealthCare Lawyer, Nov.\ 2026}); (2) \textbf{State-Level Wealth Taxes:} Because the federal Direct Tax Clause does not bind the States, billionaire mark-to-market and net-worth proposals shifted to statehouses such as California and New York through 2026; and (3) \textbf{State Tax Conformity Litigation:} Multinationals and amici have cited \textit{Moore}'s pass-through characterization of \S\ 965 in state-court Commerce Clause disputes over whether states may tax deemed-repatriated foreign affiliate earnings without adjusting state apportionment formulas (\href{https://www.nftc.org/wp-content/uploads/2026/11/2026-11-09-National-Foreign-Trade-Council-Amicus-Brief.pdf}{NFTC Amicus Brief, Or.\ Sup.\ Ct., Nov.\ 2026}).

\section{Pedagogical Critiques and Exam Hypotheticals}

\subsection{Pedagogical Critiques}
\begin{itemize}
  \item \textbf{Doctrinal Minimalism vs.\ Consequentialism:} Critics argue Kavanaugh's opinion ducked the central question granted on certiorari to avoid triggering a ``fiscal calamity,'' leaving lower courts with an undertheorized ``non-arbitrary attribution'' standard drawn from Due Process rather than a clear Sixteenth Amendment rule.
  \item \textbf{Formalism vs.\ Economic Substance:} The split between Barrett/Thomas (insisting on legal separation between corporation and shareholder) and Kavanaugh/Jackson (treating entity-level realization as sufficient) highlights a recurring constitutional law tension: whether Article I tax limits track formal legal entities or real-world economic control and revenue realities.
\end{itemize}

\subsection{Exam Hypotheticals}
\begin{enumerate}
  \item \textbf{Hypothetical 1 (Widely Held Domestic Pass-Through Tax):} Congress repeals the corporate income tax (\S\ 11) for publicly traded U.S.\ corporations and enacts a pass-through tax attributing each corporation's realized annual net profits pro rata to every U.S.\ shareholder owning at least one share on December 31. \textit{Analysis:} Under Kavanaugh's majority (see footnote 8 rejecting Barrett's domestic/widely held distinction), this satisfies pass-through attribution of entity-realized income, though challengers would press Due Process arbitrariness (lack of shareholder control or ability to compel dividends) and invoke Barrett's concurrence plus Thomas's dissent (4 Justices) to argue corporate form cannot be pierced for passive retail investors.
  \item \textbf{Hypothetical 2 (``Opposite Side'' Contrast --- Mark-to-Market Appreciation Tax on Individual Art and Stock Portfolios):} Congress enacts an unapportioned annual 2\% ``Billionaire Appreciation Tax'' on the year-over-year increase in fair market value of unsold real estate, fine art, and publicly traded equities held by individuals worth over \$100 million. \textit{Analysis:} Unlike \textit{Moore}, \textit{no entity has realized any income}---the tax targets pure unrealized appreciation of property held directly by individuals. Under \textit{Moore}'s explicit reservation, pass-through attribution cannot save this statute. Four Justices (Barrett, Alito, Thomas, Gorsuch) are already on record that unapportioned taxes on unrealized appreciation violate the Sixteenth Amendment and Direct Tax Clause; the Government would need all five remaining Justices (including Roberts and Kavanaugh, who stressed the Government's own concessions) either to reject a constitutional realization requirement (per Jackson) or to recharacterize the levy as an indirect excise.
\end{enumerate}

\end{document}
